\originalchapter{Cooper不稳定性与超导序}
\label{ch:bcs}
\source{8.511第22--24讲，衔接8.512第9--13讲。}
\originalsection{费米海上方的一对电子}
考虑总动量零、相反自旋的两电子，只允许$0<\xi_k=\epsilon_k-E_F<\hbar\omega_D$。约化吸引势为$-V_0$，配对波函数系数$a_k$满足$(2\xi_k-E)a_k=V_0\sum_{k'}a_{k'}$。消去$a_k$得
\begin{equation}
1=V_0\sum_k\frac1{2\xi_k-E}\simeq N_\uparrow(0)V_0\int_0^{\hbar\omega_D}\frac{\dd\xi}{2\xi+|E|}.
\end{equation}
$N_\uparrow(0)$是每自旋态密度。弱吸引时$|E|\simeq2\hbar\omega_D\exp[-2/(N_\uparrow(0)V_0)]$，任意小吸引都产生束缚。与真空三维两体问题的有限阈值不同，费米海边缘的常态密度导致对数积分发散。两体结果指出正常费米面不稳定，却未给具有有限配对密度的真正基态。
\originalsection{BCS变分与Bogoliubov准粒子}
多体试验态取$\ket\Psi=\prod_k(u_k+v_kc_{k\uparrow}^\dagger c_{-k\downarrow}^\dagger)\ket0$，$|u_k|^2+|v_k|^2=1$。不同$k$的空态与双占据态叠加，相位一致的异常平均$\avg{c_{-k\downarrow}c_{k\uparrow}}=u_k^*v_k$成为序参量。以$K=H-\mu N$变分，令$\Delta=V_0\sum_ku_k^*v_k$，得到$u_kv_k=\Delta/(2E_k)$、$v_k^2=(1-\xi_k/E_k)/2$，其中$E_k=\sqrt{\xi_k^2+|\Delta|^2}$。

同样结果可从Nambu矩阵获得。对$(c_{k\uparrow},c_{-k\downarrow}^\dagger)^T$，单个配对块为
\begin{equation}
\mathcal H_k=\begin{pmatrix}\xi_k&-\Delta\\-\Delta^*&-\xi_k\end{pmatrix},\qquad \mathcal H_k^2=E_k^2I.
\end{equation}
归一化本征矢形成规范变换$\gamma_{k\uparrow}=u_kc_{k\uparrow}-v_kc_{-k\downarrow}^\dagger$，仍满足费米反对易关系。准粒子是电子与空穴的相干混合，创建它至少需能量$|\Delta|$，而打破一对并产生两个准粒子通常需$2|\Delta|$。

有限温度下每个准粒子态占据$f(E_k)$，异常平均乘$1-2f(E_k)=\tanh(E_k/2\kb T)$，故
\begin{equation}
1=N_\uparrow(0)V_0\int_0^{\hbar\omega_D}\frac{\tanh(\sqrt{\xi^2+\Delta^2}/2\kb T)}{\sqrt{\xi^2+\Delta^2}}\dd\xi.
\end{equation}
$T=0$积分为$\operatorname{arsinh}(\hbar\omega_D/\Delta_0)$，弱耦合给$\Delta_0=2\hbar\omega_D\e^{-1/[N_\uparrow(0)V_0]}$。令$\Delta\to0$时积分的有限常数给$\kb T_c=(2\e^\gamma/\pi)\hbar\omega_D\e^{-1/[N_\uparrow(0)V_0]}$，因此$2\Delta_0/(\kb T_c)=2\pi\e^{-\gamma}\simeq3.53$。这是各向同性弱耦合$s$波结论；不同配对对称、强耦合与多带材料可偏离。

正常态密度近似常数时，超导态密度为$D_s(E)=D_n(0)\operatorname{Re}[|E|/\sqrt{E^2-\Delta^2}]$。平方根边缘来自把$\xi$映到$E$的Jacobian。低温热激发按$\e^{-\Delta/\kb T}$减少，节点超导则不能沿用此指数律。
\originalsection{Ginzburg--Landau、London与涡旋}
接近$T_c$，用缓慢变化的复序参量$\Psi$写
\begin{equation}
F=\int\dd^3r\left[\alpha|\Psi|^2+\frac\beta2|\Psi|^4+\frac1{2m_p}|(-\ii\hbar\nabla-q_p\vct A)\Psi|^2+\frac{B^2}{2\mu_0}\right],\quad q_p=-2e.
\end{equation}
稳定性要求$\beta>0$，$\alpha<0$时均匀极小$|\Psi|^2=-\alpha/\beta$。变分对$\Psi^*$给GL方程，对$\vct A$给电流$\vct j=(q_p\hbar/m_p)\operatorname{Im}(\Psi^*\nabla\Psi)-(q_p^2/m_p)|\Psi|^2\vct A$。取$\Psi=|\Psi|\e^{\ii\theta}$并固定幅度，得到$\vct j=(q_p|\Psi|^2/m_p)(\hbar\nabla\theta-q_p\vct A)$。在无涡旋区域取旋度，再用Maxwell方程，$\nabla^2\vct B=\vct B/\lambda^2$，$\lambda^2=m_p/(\mu_0q_p^2|\Psi|^2)$。这是Meissner排磁的空间尺度。

幅度恢复尺度为$\xi=\hbar/\sqrt{2m_p|\alpha|}$。沿远离涡核的闭合回路电流可忽略，单值性要求$\oint\nabla\theta\cdot\dd\vct l=2\pi n$，故通量绝对值$|\Phi|=n h/(2e)$。相位 winding 不能被单值规范变换消掉，涡核内幅度需减小，否则梯度能发散。$\kappa=\lambda/\xi$控制界面能：GL极限$\kappa<1/\sqrt2$为I类，$\kappa>1/\sqrt2$为II类，后者在两个临界场之间形成涡旋态。临界场近似$B_{c2}=\Phi_0/(2\pi\xi^2)$、$B_{c1}\simeq\Phi_0[\log\kappa+O(1)]/(4\pi\lambda^2)$仅适于相应连续理论与大$\kappa$范围。
\originalsection{隧穿与Josephson关系}
两个超导体的配对振幅弱耦合，最低相位依赖自由能为$F_J=-E_J\cos\varphi$，$\varphi$是扣除跨结规范线积分的相位差。相位与转移的对数互为共轭，故$I=(2e/\hbar)\partial_\varphi F_J=I_c\sin\varphi$，$I_c=2eE_J/\hbar$。电压$V$使两侧化学势相差$2eV$，$\dot\varphi=2eV/\hbar$，恒压产生频率$2eV/h$的交流超流。正常准粒子隧穿则为$\int D_1(E)D_2(E+eV)[f(E)-f(E+eV)]\dd E$的谱卷积，与零电压相干对隧穿是不同通道。
